\documentclass[a4paper]{article}
\usepackage{amsmath}
\usepackage{amssymb}
\usepackage{geometry}
\usepackage{hyperref}

\geometry{margin=25mm}

\title{Gravitational Lensing by Torsion Stars:\\
The Return of a Captured Ray}

\author{https://github.com/hypernumbernet}
\date{\today}

\begin{document}

\maketitle

\begin{abstract}
A torsion star keeps the Schwarzschild exterior, and therefore keeps the photon sphere. It does not keep the capture. The admissible ceiling freezes the redshift at a positive floor before the metric coefficient can change sign. On the three-axis reading \(1<r_\star/r_s<1.005\), and on the one-axis saturation \(r_\star/r_s=1/(1-\mathrm{e}^{-\pi})\); both lie inside \(r_{\mathrm{ph}}=\frac32 r_s\). A ray with \(b>b_c=3\sqrt{3}\,GM/c^2\) turns outside that sphere and never meets a node. A ray with \(b<b_c\) is the ray a black hole captures. On the classical chart it has no turning point, the Killing time down to \(r_s\) diverges, and for \(r<r_s\) the radial direction has become timelike. On the frozen chart the same ray turns at \(r=b\sqrt{A_{\min}}\), inside the freeze, and returns. The round trip from the photon sphere to the freeze is \(11.44\,GM/c^3\) or \(20.96\,GM/c^3\), on either side of one photon half-orbit \(3\pi\sqrt{3}\,GM/c^3\). The time spent inside the freeze is longer. Its minimum already exceeds that half-orbit: more than \(20\,GM/c^3\) on the one-axis saturation and more than \(280\,GM/c^3\) on the three-axis ceiling. A radial ray, which reaches the centre, takes more than \(80\) and more than \(400\). The equal-scale nodes lie on this finite interior. Read as thin signed shells, an attractive shell adds an interior radial critical curve once \(R>b_m/\sqrt{2}\), and an interior tangential curve only for \(b_m/\sqrt{2}<R<b_m\), outside the radial one. A repulsive shell adds neither curve on its own disk. The masses are unfixed, so magnifications are not predicted. Solar light, galactic rings, and the emission rings of M87* and Sgr A* stay outside \(r_\star\). A late echo train is not this return.
\end{abstract}

\section{Introduction}
\label{sec:introduction}

The double spacetime theory refuses a shared continuum. Each massive particle carries a usual Minkowski frame and a dual frame, and gravity is the torsional mismatch of that pair~\cite{dst-main}. On the exterior of a spherical mass the mismatch is a radial boost, and the boost reproduces the Schwarzschild coframe. Inside a sufficient density the same mismatch changes sign with radius. The parent theory calls the resulting object a torsion star: a layered condensate with no event horizon and no central singularity. The layers are the zeros of one interference factor, one node in each interval of \(\pi\), not an interior solution pasted onto the Schwarzschild chart.

Light asks a different question from the fall of a massive particle. The written action gives a photon vanishing stiffness, so it does not accelerate a photon's rapidities. What deflects light on the exterior is the null geometry of the coframe. The question of this note is what that geometry does to the rays a black hole would capture.

The Schwarzschild lens has one strong-field critical impact parameter, the photon sphere, and one weak-field critical curve, the Einstein ring of the total mass. Both survive, because both are fixed outside the freeze. The rays that pass the photon sphere do not. On the classical chart they have nowhere to turn: the Killing time down to the horizon diverges, and inside the horizon the radial coordinate has become timelike. The frozen chart gives those rays a periapsis and a finite return. That return, not a second copy of the Einstein ring, is the interior signature of the torsion star.

Two readings have to be kept apart. The exterior deflection, the place of the photon sphere relative to the quasi-horizon, the nodal window, and the return of a subcritical ray are readings of the chart. The identification of a node with a thin signed shell is a model of the barriers in the weak deflection. It does not assign a mass to each barrier, and it is not a null integration through a strong-field interior.

Section~\ref{sec:chart} records the chart. Section~\ref{sec:capture} says what the Schwarzschild geometry does with a ray inside \(b_c\). Section~\ref{sec:return} says what the freeze does instead. Section~\ref{sec:shells} reads the nodes as a weak lens. Section~\ref{sec:observations} compares the exterior class with the solar delay, with galactic rings, and with the emission rings of M87* and Sgr A*. Section~\ref{sec:discriminants} isolates the interval on which the two theories can actually be told apart.

\section{The chart}
\label{sec:chart}

On the exterior the radial boost and the redshift are
\begin{equation}
\label{eq:exterior-A}
\varphi(r)=-\frac12\log\bigl(1-r_s/r\bigr),
\qquad
A(r)=1-\frac{r_s}{r},
\qquad
r_s=\frac{2GM}{c^2},
\end{equation}
with \(r>r_s\). The coframe \(\bigl(\sqrt{A},\,A^{-1/2},\,r,\,r\sin\theta\bigr)\) induces the Schwarzschild metric. Three quadratic quantities live on this chart and are not the same object: the finite-angle mismatch \(J=\frac12\varphi^2\), the field seed \(J_{\mathrm{field}}=\frac12(\varphi')^2\), and the teleparallel scalar \(T\). They agree at leading order in the weak field and nowhere else. The Einstein equation belongs to \(T\) after a tetrad variation. It is not an identity for \(J\).

Classical Schwarzschild sends \(\varphi\to\infty\) and \(A\to 0\) as \(r\to r_s^+\), and for \(r<r_s\) the factor \(A\) is negative. The admissible cone cuts the divergence before that sign change. Reading the three-axis ceiling \(|J|=3\pi^2/8\) as a radial pure boost inverts to
\begin{equation}
\label{eq:quasihorizon}
\varphi_{\max}=\frac{\pi\sqrt{3}}{2},
\qquad
A_{\min}=\mathrm{e}^{-\pi\sqrt{3}},
\qquad
\frac{r_\star}{r_s}=\frac{1}{1-A_{\min}},
\end{equation}
with \(0<A_{\min}<10^{-2}\) and \(1<r_\star/r_s<1.005\). For \(r\le r_\star\) the chart freezes \(\varphi=\varphi_{\max}\) and \(A=A_{\min}\). The frozen factor is strictly positive, so a null surface \(A=0\) does not form, and the static Killing field remains timelike through the interior. The same passage of the parent theory flags the identification as a three-axis ceiling read onto one boost axis. A one-axis admissible boost saturates at \(\varphi=\pi/2\), where \(A=\mathrm{e}^{-\pi}\) and \(r_\star/r_s=1/(1-\mathrm{e}^{-\pi})\), still less than \(3/2\).

The layered interior is organized by a different scalar. Under the equal-scale ansatz the interference factor, in \(x=\ell/(2r)\), is
\begin{equation}
\label{eq:gamma-s}
\gamma_s(x)=\cosh x\cos x-\sinh x\sin x,
\end{equation}
with derivative \(\gamma_s'(x)=-2\cosh x\,\sin x\). Each interval \([n\pi,(n+1)\pi]\) contains exactly one node, the sign of \(\gamma_s\) alternates from plateau to plateau, and the node lies in \(n\pi+\pi/4<x_n<n\pi+\pi/2\). The first two areal radii therefore satisfy
\begin{equation}
\label{eq:radius-ratio}
\frac16<\frac{r_2}{r_1}<\frac{4}{5\pi}\approx 0.255.
\end{equation}
The ratio depends only on the dimensionless roots. The length \(\ell\) sets the absolute scale and cancels. Plateaux are screened by \(1/\cosh\bigl(\ell/(2r)\bigr)\). What grows toward the centre is the barrier at the node. On the admissible cone \(|J|\le 3\pi^2/8\), and the teleparallel density stays in \(-8<r^2 T<27\), so the spike is finite. The cone does not convert that window into a unique surface density.

A photon, in the written action, is a particle of vanishing stiffness. Vanishing stiffness removes the acceleration of the common rapidity. It does not force the photon onto the vacuum, and it does not make \(F=k/(\gamma_s r^2)\) the force on a ray. Null deflection on the exterior is the geodesic of the coframe~\eqref{eq:exterior-A}. Null deflection through the barriers is the model of Section~\ref{sec:shells}.

\section{Capture}
\label{sec:capture}

A static spherical line element
\[
ds^2=-A(r)\,c^2\,dt^2+A(r)^{-1}\,dr^2+r^2\,d\Omega^2
\]
conserves an energy \(E\) and an angular momentum \(L\). Their ratio is the asymptotic impact parameter \(b=L/E\). A radial turning point would be a radius at which
\begin{equation}
\label{eq:turning}
b^2=\frac{r^2}{A(r)}.
\end{equation}
For the exterior factor the right-hand side, written in \(u=r/r_s>1\), is \(r_s^2\,u^3/(u-1)\). The cubic identity
\[
4u^3-27u+27=(2u-3)^2(u+3)
\]
says that this quantity is at least \(\frac{27}{4}r_s^2\), with equality only at \(u=\frac32\). That minimum is the photon sphere
\begin{equation}
\label{eq:photon-sphere}
r_{\mathrm{ph}}=\frac32 r_s=\frac{3GM}{c^2},
\qquad
b_c=\frac{r_{\mathrm{ph}}}{\sqrt{A(r_{\mathrm{ph}})}}=3\sqrt{3}\,\frac{GM}{c^2}.
\end{equation}
The circular orbit is unstable because the minimum of \(r^2/A\) is a maximum of the effective potential \(A/r^2\). A ray from infinity with \(b>b_c\) has a turning point at some \(r_0>r_{\mathrm{ph}}\). A ray with \(b<b_c\) has none on \(r>r_s\): every classical radius demands a larger impact parameter than the ray carries, so the ray does not turn.

The same chart makes the inward trip infinitely long. A radial null ray satisfies \(c\,dt=dr/A\), and
\begin{equation}
\label{eq:antideriv}
\frac{d}{dr}\Bigl(r+r_s\log|r-r_s|\Bigr)=\frac{1}{A(r)}
\end{equation}
on \(r>r_s\). As \(r\downarrow r_s\) the logarithm diverges, so for every proposed delay there is a radius, still outside the horizon, at which the Killing time from the photon sphere already exceeds it. The ray reaches \(r_s\) only at infinite Killing time, and it does not come back.

Inside the horizon the obstruction is sharper than an infinite delay. For \(0<r<r_s\) one has \(A<0\). The coefficient of \(dt^2\) and the coefficient of \(dr^2\) have exchanged signs: \(t\) is spacelike and \(r\) is timelike. A turning point of positive impact parameter would require \(b^2=r^2/A\) with a negative right-hand side, which is impossible. The ray is not free to reverse its radial motion. That is the interior of the black hole as the classical chart describes it.

The azimuthal advance of a ray that does turn, from a turning point \(r_0>r_{\mathrm{ph}}\) out to infinity and back, is
\begin{equation}
\label{eq:deflection-integral}
\Delta\varphi(b)
  =2\int_{r_0}^{\infty}
    \frac{dr}{r^2\sqrt{b^{-2}-A(r)/r^2}},
\qquad
\alpha_{\mathrm{ext}}(b)=\Delta\varphi(b)-\pi.
\end{equation}
In the weak field it reduces to \(4GM/(c^2 b)\). The same factor fixes the Shapiro delay. In the parametrized post-Newtonian writing both scale with \((1+\gamma)/2\), and the chart~\eqref{eq:exterior-A} is the case \(\gamma=1\).

For a thin lens the lens equation relates the angular source position \(\beta\) to the image position \(\theta\) by \(\beta=\theta-\bar\alpha(\theta)\), with \(\bar\alpha=(D_{LS}/D_S)\,\alpha(D_L\theta)\). The Einstein angle of a point mass \(M\) is
\begin{equation}
\label{eq:einstein}
\theta_E^2=\frac{4GM}{c^2}\,\frac{D_{LS}}{D_L D_S},
\end{equation}
so that the weak deflection is \(\bar\alpha=\theta_E^2/\theta\). The tangential eigenvalue of the lens map vanishes where \(\bar\alpha=\theta\), and the radial eigenvalue where \(d\bar\alpha/d\theta=1\)~\cite{NarayanBartelmann1996}. A point mass has the single positive root \(\theta=\theta_E\) and no radial root. That is the weak-field lens a torsion star must exceed if its layers are to draw further foci. It is not the distinction between capture and return.

\section{Return}
\label{sec:return}

Join the exterior factor~\eqref{eq:exterior-A} for \(r>r_\star\) to the frozen value \(A=A_{\min}\) for \(r\le r_\star\). The join is continuous. On the frozen interval \(A/r^2\) is strictly decreasing in \(r\), so the freeze itself contributes no photon sphere. Both ceilings lie inside \(r_{\mathrm{ph}}\): the three-axis ratio is less than \(1.005\), and the one-axis ratio \(1/(1-\mathrm{e}^{-\pi})\) is less than \(3/2\). The only barrier a ray from infinity must clear remains the Schwarzschild photon sphere, at the impact parameter~\eqref{eq:photon-sphere}. Layers inside \(r_\star\) do not move it.

The rays therefore split into the same two classes as in the classical problem, and then differ.

A ray with \(b>b_c\) turns at \(r_0>r_{\mathrm{ph}}>r_\star\). The whole of the deflected path is in the Schwarzschild domain. No node is encountered. Every weak-field test whose impact parameter is large compared with \(r_s\), and every strong-field test whose observable is the angular diameter \(2b_c/D_A\) of this critical curve, belongs to this class.

A ray with \(0<b<b_c\) is the ray Section~\ref{sec:capture} declared captured. On the joined chart the declaration fails. There is still no turning point on \(r>r_s\), so the ray passes the photon sphere and enters the freeze. Inside the freeze \(A\) is the positive constant \(A_{\min}\), and~\eqref{eq:turning} has the unique positive root
\begin{equation}
\label{eq:periapsis}
r_{\mathrm{peri}}=b\sqrt{A_{\min}}.
\end{equation}
For both ceilings this periapsis lies inside \(r_\star\), even in the limit \(b\uparrow b_c\). The ray that has only just cleared the photon sphere still turns inside the freeze, not in the cavity between \(r_\star\) and \(r_{\mathrm{ph}}\). Because \(A\) has not changed sign, the radial direction is still spacelike, the static observers of the chart still exist, and the ray reverses. There is no horizon to absorb it. It can recross \(r_{\mathrm{ph}}\) and reach infinity. Near \(b_c\) the passage past the unstable orbit can wind many times before the fall. That winding is the photon-ring sequence of the exterior geometry, one critical curve visited repeatedly. It is not a sequence of nodal limbs.

The Killing time of the trip splits at \(r_\star\). From the photon sphere down to the freeze and back,~\eqref{eq:antideriv} gives
\begin{equation}
\label{eq:echo-delay}
\Delta t
  =\frac{2}{c}\int_{r_\star}^{r_{\mathrm{ph}}}\frac{dr}{A(r)}
  =\frac{2}{c}\Bigl[r+r_s\log|r-r_s|\Bigr]_{r_\star}^{r_{\mathrm{ph}}}.
\end{equation}
The integral is finite once \(r_\star>r_s\). For \(A_{\min}=\mathrm{e}^{-\pi\sqrt{3}}\) it equals \(20.96\,GM/c^3\). For the one-axis saturation \(A=\mathrm{e}^{-\pi}\) it equals \(11.44\,GM/c^3\). The parent theory flags the three-axis reading as a ceiling transplanted onto one axis, so until that identification is settled the chart predicts an interval in \(\{11.44,\,20.96\}\,GM/c^3\), not a single number.

The natural clock already present outside is the half-period of the unstable photon orbit. On \(r=3GM/c^2\) the coordinate angular velocity is \(d\varphi/dt=c^3/(3\sqrt{3}\,GM)\), so
\begin{equation}
\label{eq:half-orbit}
T_{1/2}=3\pi\sqrt{3}\,\frac{GM}{c^3},
\end{equation}
which lies between \(16\) and \(17\) in these units, numerically \(16.32\,GM/c^3\). The two cavity returns fall on either side of it, by about \(4.9\,GM/c^3\) and \(4.6\,GM/c^3\).

Inside the freeze the integrand is the constant \(1/A_{\min}\). A radial ray, whose periapsis is the centre, spends
\begin{equation}
\label{eq:interior-radial}
\Delta t_{\mathrm{rad}}=\frac{2r_\star}{c\,A_{\min}}
  =\frac{4}{A_{\min}(1-A_{\min})}\,\frac{GM}{c^3}
\end{equation}
on the round trip from \(r_\star\) to the centre and back. Because \(A_{\min}<1/20\) on the one-axis saturation this already exceeds \(80\,GM/c^3\), and because \(A_{\min}<1/100\) on the three-axis ceiling it exceeds \(400\,GM/c^3\). The ray that only just passes the photon sphere turns earlier, at \(r\) approaching \(b_c\sqrt{A_{\min}}\) from below, and its interior round trip is shorter. It is still longer than one photon half-orbit: more than \(20\,GM/c^3\) on the one-axis saturation and more than \(280\,GM/c^3\) on the three-axis ceiling. The floors are not sharp. They are large enough to put every returning ray, however close to \(b_c\), on the far side of the photon clock. The classical chart has no such finite interior.

The nodes of Section~\ref{sec:chart} sit on this finite interval. Infinitely many of them accumulate toward the centre, and the coordinate time across the whole interval remains finite because the frozen speed \(dr/dt=c A_{\min}\) does not itself go to zero. The nodes are barriers on a return, not a second horizon.

\section{The nodes as a weak lens}
\label{sec:shells}

The plateaux of \(\gamma_s\) are exponentially quiet, so the smooth column they contribute is neglected. Each node is a barrier. The model of this section collapses the finite spike allowed by the cone to a thin spherical shell of areal radius \(R\) and mass \(m\), which may be negative: an attractive barrier or a repulsive one. Nothing in the cone fixes \(m\).

A line of sight at impact parameter \(b<R\) meets the shell twice. The projected column diverges at the limb in the ideal thin limit, as \(2\sigma R/\sqrt{R^2-b^2}\). The mass inside the cylinder does not. For \(b\le R\),
\begin{equation}
\label{eq:mcyl}
M_{\mathrm{cyl}}(b)=\frac{m}{R}\Bigl(R-\sqrt{R^2-b^2}\Bigr),
\end{equation}
and \(M_{\mathrm{cyl}}(b)=m\) for \(b\ge R\). The weak-field deflection, with the relativistic factor \(4G/c^2\), is \(\alpha(b)=4G\,M_{\mathrm{cyl}}(b)/(c^2 b)\). Outside the shell this is the point-mass law. On the axis \(\alpha\to 0\), as spherical symmetry requires. The slope of \(M_{\mathrm{cyl}}\) diverges as \(b\to R^-\).

Let \(b_m\) be the Einstein radius of the mass on the shell and set \(c=(b_m/R)^2\), negative when that mass is negative. In units of the shell's angular radius, the tangential ratio and the radial slope inside the shell are
\begin{equation}
\label{eq:shell-profiles}
\frac{\bar\alpha}{\theta}=\frac{c}{1+\sqrt{1-\xi^2}},
\qquad
\frac{d\bar\alpha}{d\theta}
  =\frac{c}{\sqrt{1-\xi^2}\,\bigl(1+\sqrt{1-\xi^2}\bigr)}.
\end{equation}
On the axis both equal \(c/2\). The tangential ratio only doubles on the way out, and it meets the finite value \(c\) at the limb. The radial slope does not stay finite: it runs to \(+\infty\) for an attractive shell and to \(-\infty\) for a repulsive one. What diverges at the limb is the radial eigenvalue.

For an attractive shell the two comparisons fall into a window measured by \(b_m\). There is a unique interior radial critical curve precisely when \(R>b_m/\sqrt{2}\), and a unique interior tangential critical curve precisely when
\begin{equation}
\label{eq:shell-window}
\frac{b_m}{\sqrt{2}}<R<b_m.
\end{equation}
In that window the tangential curve lies outside the radial one. Outside the shell the map is the point mass of the shell's own mass: an Einstein ring when \(R<b_m\), and never a radial critical curve. A repulsive shell keeps both the tangential ratio and the interior slope negative, so its disk contains neither curve.

The mass on the shell is unfixed, so the window is a criterion and not a count of images. The nodal ratio~\eqref{eq:radius-ratio} remains the ratio of the nodes. Neighboring critical curves are not obliged to stand in it. Smoothing the limb to the finite width required by the cone replaces the infinite radial slope by a steep peak; it does not move the tangential curve onto the limb. That curve remains where the enclosed-mass ratio equals one.

A stack changes the accounting by a finite interior mass, not by cancelling the radial divergence of the shell being crossed. The sign of \(\gamma_s\) alternates from node to node. The type of the critical curve does not have to alternate with it. The photon-ring sequence is a third arrangement entirely: successive half-orbits of the single curve \(b_c\), demagnified by \(\mathrm{e}^{-\pi}\) per half-orbit and converging to that one curve~\cite{Johnson2020}. An exponential pile-up at \(b_c\), a pair of nodes in the ratio~\eqref{eq:radius-ratio}, and a pair of thin-shell critical curves fixed by~\eqref{eq:shell-window} are three arrangements of light on the sky. This note does not integrate the ray that both winds at \(r_{\mathrm{ph}}\) and crosses the barriers.

\section{What is already seen}
\label{sec:observations}

The solar tests sit in the exterior class. A ray that grazes the photosphere has an impact parameter larger than \(r_s\) by a factor of order \(10^5\), and it does not enter a nucleus. The deflection is the weak-field law at the solar limb, and the delay is the Shapiro delay of the same factor \(A\). Cassini tracking gives
\begin{equation}
\label{eq:cassini}
\gamma=1+(2.1\pm 2.3)\times 10^{-5}
\end{equation}
\cite{Bertotti2003}. The chart is the case \(\gamma=1\). The measurement constrains a correction that would rescale the exterior bending and delay together. It does not constrain an interior the ray does not meet.

Galactic strong lensing is the same class in a weaker field. The Einstein angle and the angular size of the quasi-horizon are related by \(\theta_E^2\sim(r_s/D_L)(D_{LS}/D_S)\), so their ratio is the square root of an angle in radians, many orders of magnitude below unity for any galactic lens. Structure at \(r_\star\) does not move the Einstein ring, the image positions, or the time delay at the precision of those observations. A double Einstein ring formed by two sources at two redshifts is two exterior-class rings. Flux-ratio anomalies among images of a single quasar are likewise an exterior-class question. This note does not assign them to torsional nodes.

The emission ring of M87* has diameter \(42\pm 3\,\mu\mathrm{as}\), surrounding a central depression of flux ratio \(\gtrsim 10{:}1\). The angular gravitational radius adopted with that ring is \(\theta_g=3.8\pm 0.4\,\mu\mathrm{as}\), and the critical curve~\eqref{eq:photon-sphere} has angular diameter \(6\sqrt{3}\,\theta_g\)~\cite{EHT2019}. Since \(6\sqrt{3}\approx 10.392\), that diameter is about \(39.5\pm 4.2\,\mu\mathrm{as}\), which overlaps \(42\pm 3\,\mu\mathrm{as}\). The published ring is lensed emission around the critical curve, not the curve itself, and the mass inference already uses a shadow scaling of that order. The overlap checks that the photon sphere lies where the exterior chart places it. It is not a detection of a horizon. The horizon is what the quasi-horizon replaces, and the critical curve does not consult the replacement.

Sgr A* is the same comparison at a second mass. The emission ring has diameter \(51.8\pm 2.3\,\mu\mathrm{as}\)~\cite{EHT2022I}. Calibrated to general-relativistic simulations, the angular gravitational radius is \(4.8^{+1.4}_{-0.7}\,\mu\mathrm{as}\)~\cite{EHT2022IV}, and \(6\sqrt{3}\) times that interval runs from about \(43\) to \(64\,\mu\mathrm{as}\). The measured ring lies inside it. Both images are single rings. What the agreement leaves undecided is the return of Section~\ref{sec:return}.

\section{Where the predictions differ}
\label{sec:discriminants}

The agreements of Section~\ref{sec:observations} do not separate the joined chart from the Schwarzschild lens. The rays that produce them stay outside \(r_\star\), so \(2b_c/D_A\) is the Schwarzschild diameter, and the Cassini delay and the galactic rings have impact parameters far outside \(r_s\). On Earth baselines the next-generation array expects to separate at most the \(n=0\) and \(n=1\) subrings of the single curve \(b_c\), which already requires super-resolution~\cite{Johnson2023}. That separation does not measure a ratio bounded by \(4/(5\pi)\).

The brightness inside the shadow does not separate them either. A penetrating ray returns, so the interior need not be empty, and an optically thick plasma can keep it dark. The M87* contrast \(\gtrsim 10{:}1\) limits the sum. A surface that thermalizes the incident power, or that reflects it perfectly, is a further object, and the Sgr A* spectrum and image disfavor those two~\cite{EHT2022VI}. The freeze \(A=A_{\min}\) is neither of them.

What the joined chart adds is a finite Killing time for a ray the classical chart does not return. The cavity piece~\eqref{eq:echo-delay} is one ringdown cycle, \(11.44\,GM/c^3\) or \(20.96\,GM/c^3\), set beside the photon half-orbit~\eqref{eq:half-orbit}. The interior piece is larger. Every ray that passes the photon sphere spends more than \(20\,GM/c^3\) inside the one-axis freeze, and more than \(280\,GM/c^3\) inside the three-axis freeze, before it can start back; a radial ray spends more than \(80\) and more than \(400\). For a mass \(4\times 10^6\,M_\odot\), where \(GM/c^3\approx 20\,\mathrm{s}\), the half-orbit is a few minutes and the interior floors are longer. For the M87* mass of Section~\ref{sec:observations} they are days. A flare timed to a few \(GM/c^3\) could in principle separate one return from one photon half-orbit. No such separation is in hand.

The same interval shows why a late gravitational-wave echo is not the prediction. Echo searches bound reflecting surfaces whose proper distance to the horizon is tiny, and they look for a pulse train long after the ringdown. The round trip computed here is one cycle of that ringdown, or a few, not that train. A non-detection of a late train is what a cavity as wide as \(r_\star-r_s\), and an interior crossed at speed \(c A_{\min}\), lead one to expect. It does not confirm the cavity. The chart supplies the metric coefficient \(A\). It does not supply a boundary condition for a gravitational wave, and the shell masses remain unfixed.

Periapsis limbs remain a weaker target. They are accessible from infinity only for areal radii below \(b_c\sqrt{A_{\min}}\). Nodes outside that cut are crossed by every penetrating ray and do not appear as separate grazes. Nodes inside it would be grazed in the areal ratio~\eqref{eq:radius-ratio}, after the ceiling is chosen and only if the limb is bright enough to see.

\section{Conclusion}
\label{sec:conclusion}

The exterior of a torsion star lenses light as the Schwarzschild chart lenses light. The photon sphere lies outside both recorded ceilings, the critical impact parameter remains \(b_c=3\sqrt{3}\,GM/c^2\), and the Einstein ring of the total mass survives in the weak field. That is why the solar delay, the galactic rings, and the emission rings of M87* and Sgr A* meet the theory at the precision those measurements actually have.

The interior does not meet the black hole. A ray with \(b<b_c\) has no turning point on the classical exterior, the Killing time down to \(r_s\) is infinite, and inside \(r_s\) the metric coefficient has changed sign. The freeze keeps the coefficient positive. The same ray turns at \(r=b\sqrt{A_{\min}}\), inside \(r_\star\), and the Killing time of even the shortest such visit exceeds one photon half-orbit. The nodes lie on that finite return. Read as thin signed shells they can add critical curves, an interior radial curve once \(R>b_m/\sqrt{2}\) and an interior tangential curve only inside the window \(b_m/\sqrt{2}<R<b_m\), while a repulsive shell adds neither on its own disk. The mass on each shell is not determined, so the model does not predict which window is occupied, nor the magnifications, nor the deflection of a ray that both winds at the photon sphere and crosses the barriers.

The interval the chart itself adds is therefore a time, not a second ring. General relativity does not return the ray. The torsion star returns it from inside the freeze, on a clock that is not the photon half-orbit and not a late echo train.

\begin{thebibliography}{99}
\raggedright

\bibitem{dst-main}
Anonymous,
``Gravity as Torsion between Dual Spacetime: A Biquaternionic Reformulation of General Relativity'' (2025).
\url{https://github.com/hypernumbernet/dual-spacetime-doc/blob/main/double-spacetime-theory.pdf}.

\bibitem{NarayanBartelmann1996}
R.~Narayan and M.~Bartelmann,
``Lectures on Gravitational Lensing,''
arXiv:astro-ph/9606001 (1996).

\bibitem{Bertotti2003}
B.~Bertotti, L.~Iess, and P.~Tortora,
``A test of general relativity using radio links with the Cassini spacecraft,''
\textit{Nature} \textbf{425}, 374--376 (2003).

\bibitem{EHT2019}
Event Horizon Telescope Collaboration,
``First M87 Event Horizon Telescope Results. I. The Shadow of the Supermassive Black Hole,''
\textit{Astrophys.\ J.\ Lett.} \textbf{875}, L1 (2019).

\bibitem{EHT2022I}
Event Horizon Telescope Collaboration,
``First Sagittarius A* Event Horizon Telescope Results. I. The Shadow of the Supermassive Black Hole in the Center of the Milky Way,''
\textit{Astrophys.\ J.\ Lett.} \textbf{930}, L12 (2022).

\bibitem{EHT2022IV}
Event Horizon Telescope Collaboration,
``First Sagittarius A* Event Horizon Telescope Results. IV. Variability, Morphology, and Black Hole Mass,''
\textit{Astrophys.\ J.\ Lett.} \textbf{930}, L15 (2022).

\bibitem{EHT2022VI}
Event Horizon Telescope Collaboration,
``First Sagittarius A* Event Horizon Telescope Results. VI. Testing the Black Hole Metric,''
\textit{Astrophys.\ J.\ Lett.} \textbf{930}, L17 (2022).

\bibitem{Johnson2020}
M.~D.~Johnson et al.,
``Universal interferometric signatures of a black hole's photon ring,''
\textit{Sci.\ Adv.} \textbf{6}, eaaz1310 (2020).

\bibitem{Johnson2023}
M.~D.~Johnson et al.,
``Key Science Goals for the Next-Generation Event Horizon Telescope,''
\textit{Galaxies} \textbf{11}, 61 (2023).

\end{thebibliography}

\end{document}
